\chapter{Anatomia de um poço submarino}
\label{cap:dominio}

\begin{objetivos}
Ao final deste capítulo, você será capaz de:
\begin{itemize}[leftmargin=*,itemsep=0.2em]
  \item descrever os componentes de um poço submarino de produção, do reservatório à
        unidade flutuante, e nomear as válvulas que compõem a árvore de natal molhada;
  \item explicar a função dos sensores PDG e TPT e justificar por que exatamente essas
        quatro variáveis dominam os métodos deste livro;
  \item caracterizar cada uma das nove classes de anomalia do conjunto 3W em termos de
        mecanismo físico e assinatura esperada nos sensores;
  \item distinguir eventos de causa-raiz de eventos sintomáticos, e explicar por que
        essa distinção afeta o projeto do classificador;
  \item descrever a completação inteligente e os eventos associados a válvulas de
        controle de intervalo (ICV).
\end{itemize}
\end{objetivos}

\section{Do reservatório à plataforma}

Um sistema de produção \ing{offshore} pode ser lido como uma sequência de quatro
domínios, cada um com sua física e seus modos de falha.

\begin{enumerate}[leftmargin=*,itemsep=0.5em]
  \item \textbf{Reservatório.} A rocha porosa que contém os hidrocarbonetos. Fornece a
        energia primária do escoamento (pressão estática) e determina a composição do
        fluido produzido --- razão gás-óleo, fração de água, presença de
        CO\textsubscript{2} e H\textsubscript{2}S.
  \item \textbf{Poço.} A conexão vertical (ou desviada) entre reservatório e leito
        marinho, revestida por colunas de aço cimentadas e equipada internamente com a
        \termo{coluna de produção}{production tubing} e seus acessórios.
  \item \textbf{Sistema submarino.} A árvore de natal molhada instalada na cabeça do
        poço, o \ing{jumper} que a conecta ao \ing{manifold} ou diretamente ao
        \ing{riser}, e as linhas de fluxo que atravessam o leito marinho.
  \item \textbf{Unidade de superfície.} O \ing{riser} que sobe até a unidade flutuante,
        o separador trifásico, o sistema de compressão e o de tratamento de água.
\end{enumerate}

Neste livro, o objeto de interesse é predominantemente o segundo e o terceiro domínio.
Os eventos indesejáveis do conjunto 3W ocorrem entre o fundo do poço e a válvula
reguladora de produção, e é ali que estão os sensores que os revelam.

\subsection{A completação}

Chama-se \termo{completação}{well completion} o conjunto de operações e equipamentos
que preparam o poço para produzir de forma controlada. Os elementos relevantes para
este livro são:

\begin{description}[leftmargin=0pt,style=nextline,itemsep=0.5em]

\item[Coluna de produção.] Tubo interno por onde o fluido escoa até a superfície. Isola
o fluxo do revestimento e permite a instalação de acessórios.

\item[\ing{Packer}.] Elemento de vedação que ancora a coluna ao revestimento e isola o
espaço anular, impedindo comunicação de pressão entre a zona produtora e o restante do
poço.

\item[Válvula de segurança de subsuperfície (DHSV).] Instalada tipicamente a algumas
centenas de metros abaixo do leito marinho, é uma válvula do tipo \ing{flapper}
mantida aberta por pressão hidráulica aplicada por uma linha de controle. \emph{A
falha segura é o fechamento}: se a pressão hidráulica se perde por qualquer motivo, a
válvula fecha. Essa escolha de projeto --- correta do ponto de vista de segurança ---
é a origem do fechamento espúrio, um dos eventos mais frequentes no conjunto 3W.

\item[Válvulas de gás de elevação (GLV).] Em poços com elevação artificial por gás,
permitem a injeção de gás do anular para a coluna, reduzindo a densidade da mistura e,
portanto, a pressão hidrostática que o reservatório precisa vencer.

\item[Válvulas de controle de intervalo (ICV).] Em poços de \termo{completação
inteligente}{intelligent completion}, permitem regular ou fechar remotamente cada zona
produtora individualmente. São o assunto de uma classe de anomalia inteira
(\cref{sec:dominio-icv}) e aparecem com destaque no \cref{cap:mundoaberto}.

\end{description}

\subsection{A árvore de natal molhada}

A \termo{árvore de natal molhada}{wet christmas tree} (ANM) é o conjunto de válvulas
instalado sobre a cabeça do poço no leito marinho. Sua função é dar controle de fluxo,
permitir intervenções e garantir barreiras de segurança redundantes. A nomenclatura das
válvulas varia entre operadoras; a que adotamos aqui é a usual na literatura brasileira
e no conjunto 3W:

\begin{center}
\small
\begin{tabular}{ll p{7.5cm}}
\toprule
\textbf{Sigla} & \textbf{Nome} & \textbf{Função} \\
\midrule
M1 & Master 1        & Válvula mestre inferior da linha de produção; primeira barreira. \\
M2 & Master 2        & Válvula mestre superior; segunda barreira em série com M1. \\
W1 & Wing 1          & Válvula de asa da linha de produção. \\
W2 & Wing 2          & Válvula de asa da linha de serviço ou anular. \\
XO & \ing{Cross-Over}& Interliga produção e anular; usada em operações especiais. \\
PXO & \ing{Production Cross-Over} & Cruzamento dedicado à produção. \\
CKP & \ing{Production Choke} & Válvula reguladora de produção; controla a vazão por
    estrangulamento. Também chamada PCK. \\
CKGL & \ing{Gas-Lift Choke} & Válvula reguladora da injeção de gás de elevação. \\
\bottomrule
\end{tabular}
\end{center}

\begin{destaque}[PCK ou CKP?]
As duas siglas designam a mesma válvula reguladora de produção. O conjunto 3W nomeia
duas de suas classes com ``PCK'' (\emph{Quick Restriction in PCK} e \emph{Scaling in
PCK}), enquanto a documentação de engenharia usa com frequência ``CKP''. Neste livro
mantemos ``PCK'' ao nos referirmos às classes do 3W, para preservar a rastreabilidade,
e ``CKP'' ao discutirmos o equipamento.
\end{destaque}

\section{O que os sensores enxergam}
\label{sec:dominio-sensores}

Aqui está o fato que molda todo o restante do livro: \textbf{o número de variáveis
observadas é pequeno, e elas são fortemente correlacionadas}.

\begin{figure}[htbp]
\centering
\includegraphics[width=0.42\textwidth]{esquema_vazio.png}
\caption{Disposição típica de válvulas e sensores em um poço submarino. Os medidores
PDG, TPT e PT ocupam posições fixas ao longo do caminho do fluido, e cada válvula ---
DHSV, M1, M2, W1, W2, XO, PXO e a reguladora --- define um estado operacional distinto.
Adaptado de \citet{aranha2023system}.}
\label{fig:dominio-esquema}
\end{figure}

A \cref{fig:dominio-esquema} situa os sensores no caminho do fluido. Essa geometria
explica boa parte do que vem a seguir: o PDG está a montante de tudo e responde primeiro
a variações no reservatório; o TPT está no leito marinho, onde o resfriamento é máximo;
o PT está a jusante da reguladora, e portanto enxerga a pressão já estrangulada. Duas
anomalias de mesma origem física produzem assinaturas diferentes em cada um desses
pontos, e é essa diferença que os modelos exploram.

\subsection{Os quatro sinais fundamentais}

\begin{description}[leftmargin=0pt,style=nextline,itemsep=0.5em]

\item[P-PDG --- pressão no medidor permanente de fundo.]
O \termo{medidor permanente de fundo}{permanent downhole gauge} (PDG) é instalado na
coluna de produção, próximo à zona produtora. Mede a pressão no ponto mais próximo do
reservatório que a instrumentação alcança. É o sinal mais informativo sobre o
comportamento do reservatório e sobre eventos de fundo, e é também o mais sujeito a
falha --- sua substituição exige intervenção com sonda.

\item[T-PDG --- temperatura no medidor permanente de fundo.]
Mede a temperatura do fluido no fundo. Varia lentamente em condições normais e reage a
mudanças de vazão com constante de tempo apreciável, o que a torna útil para
discriminar transitórios rápidos de mudanças de regime.

\item[P-TPT --- pressão no transdutor de temperatura e pressão.]
O \termo{transdutor de temperatura e pressão}{temperature and pressure transducer}
(TPT) está instalado na árvore de natal, a jusante do poço e a montante da válvula
reguladora. Mede a pressão na cabeça do poço.

\item[T-TPT --- temperatura no transdutor de temperatura e pressão.]
Temperatura na cabeça do poço. É a variável mais diretamente ligada à formação de
hidratos, já que a região da ANM é onde a combinação de baixa temperatura ambiente e
alta pressão é mais severa.

\end{description}

O conjunto 3W disponibiliza outras variáveis --- pressão a jusante da \ing{choke}
(P-JUS-CKP), temperatura a jusante, vazões de gás de elevação, posição da \ing{choke}.
Contudo, a maior parte dos trabalhos discutidos neste livro utiliza o subconjunto
\{P-PDG, T-PDG, P-TPT, T-TPT\}, por três razões: são as variáveis presentes em
praticamente todas as instâncias, são as menos afetadas por convenções de instalação
específicas de cada poço, e carregam a informação física essencial dos eventos de
interesse.

\begin{exemplo}{O gradiente entre fundo e cabeça}
A diferença $P_{\text{PDG}} - P_{\text{TPT}}$ é, aproximadamente, a perda de carga ao
longo da coluna de produção --- soma da componente hidrostática com a componente
friccional. Quando o fluido fica mais denso (aumento de BSW, por exemplo), a
componente hidrostática cresce e o gradiente aumenta. Quando a vazão cai, a componente
friccional diminui e o gradiente diminui.

Um único par de sensores permite, portanto, inferir mudanças em duas propriedades
distintas do escoamento. É por essa razão que o agente de linguagem do
\cref{cap:agente} calcula explicitamente um descritor de \emph{relação entre P-PDG e
P-TPT} --- ``mesma direção'', ``divergentes'' ou ``um estável'' --- e o entrega ao
modelo como evidência. A informação está na relação, não em cada sinal isolado.
\end{exemplo}

\subsection{Por que os sensores mentem}

Três patologias de instrumentação atravessam todo o livro e precisam ser enunciadas
desde já.

\begin{itemize}[leftmargin=*,itemsep=0.4em]
  \item \textbf{Congelamento.} O sensor para de atualizar e repete indefinidamente a
        última leitura válida. Uma série congelada tem desvio-padrão nulo e é
        trivialmente detectável --- desde que alguém procure por isso.
  \item \textbf{Perda de telemetria.} O valor se torna ausente ou é substituído por
        zero. Essa patologia terá importância inesperada no
        \cref{cap:resultados-agente}: em várias classes do 3W, o início da anomalia
        coincide com a perda dos medidores de fundo, o que cria uma correlação espúria
        entre ``anomalia'' e ``ausência de dado''.
  \item \textbf{Deriva.} O sensor mantém a dinâmica, mas seu nível absoluto se desloca
        lentamente. É a mais perigosa das três, porque não produz nenhuma assinatura
        óbvia e corrompe silenciosamente qualquer modelo calibrado em valores
        absolutos.
\end{itemize}

\begin{destaque}[Consequência de projeto]
As três patologias acima são o argumento técnico para trabalhar com \emph{variações
relativas a uma linha de base} em vez de valores absolutos. Essa decisão aparece de
formas diferentes ao longo do livro: normalização por janela no \cref{cap:dual},
padronização por poço no \cref{cap:segmentacao} e, de forma mais explícita, no cálculo
do descritor $\delta = \bar{e} - \bar{b}$ do \cref{cap:agente}.
\end{destaque}

\section{A taxonomia dos eventos indesejáveis}
\label{sec:dominio-taxonomia}

O conjunto 3W organiza os eventos em uma classe normal e nove classes de anomalia
\citep{VARGAS2019,VazVargas2026_3w2}. Descrevemos cada uma pelo mecanismo físico e pela
assinatura esperada.

\subsection{Classe 0 --- Estado normal}

Operação estável. Os sinais apresentam variação de baixa amplitude em torno de níveis
determinados pelo ponto de operação. É a classe de referência contra a qual todas as
demais são contrastadas e, em quantidade de amostras, a mais numerosa por larga
margem.

\subsection{Classe 1 --- Aumento abrupto de BSW}

\textbf{Mecanismo.} O \termo{teor de água e sedimentos}{basic sediment and water} (BSW)
é a fração volumétrica de água no fluido produzido. Um aumento abrupto indica avanço
da frente de água do aquífero ou de um poço injetor, ou ainda comunicação com uma zona
aquífera por falha de isolamento.

\textbf{Assinatura.} Água é mais densa que óleo. O aumento da fração de água eleva a
massa específica média da coluna, aumentando a pressão hidrostática. Espera-se aumento
de P-PDG e mudança no gradiente entre fundo e cabeça. A temperatura pode cair, já que a
água tem capacidade calorífica maior.

\textbf{Dificuldade.} O evento é gradual em muitos casos reais, o que o confunde com
declínio natural. A classe 1 apresenta desempenho apenas mediano em quase todos os
experimentos deste livro.

\subsection{Classe 2 --- Fechamento espúrio da DHSV}

\textbf{Mecanismo.} Perda de pressão na linha de controle hidráulico leva a DHSV a
fechar sem comando. O poço é isolado abaixo da válvula.

\textbf{Assinatura.} A mais nítida de todas. A pressão a montante da válvula (P-PDG)
sobe rapidamente, pois o reservatório continua a alimentar um volume fechado; a pressão
a jusante (P-TPT) despenca. As duas pressões divergem de forma inequívoca. As
temperaturas caem, pois cessa o transporte convectivo de calor do fundo.

\textbf{Consequência para os métodos.} A classe 2 é aquela em que praticamente todo
método deste livro atinge desempenho excelente --- 95\,\% no ensemble binário do
\cref{cap:closedworld}, 100\,\% na detecção de novidade do \cref{cap:resultados-agente}.
É o caso fácil, e serve de controle positivo: um método que falhe na classe 2 tem um
defeito estrutural.

\subsection{Classe 3 --- Golfadas severas}

\textbf{Mecanismo.} Em escoamento multifásico ascendente com \ing{riser} em catenária,
o líquido pode acumular-se no ponto baixo até bloquear a passagem de gás. A pressão a
montante cresce até vencer a coluna de líquido, que é então expelida em golfada. O
ciclo se repete com período de minutos a dezenas de minutos.

\textbf{Assinatura.} Oscilação periódica de grande amplitude em P-TPT e, com menor
amplitude, em P-PDG. É a única classe cuja assinatura característica é \emph{de
frequência}, não de nível.

\textbf{Consequência para os métodos.} Métodos que resumem uma janela por média e
desvio-padrão perdem a informação de periodicidade. Esse é um dos motivos pelos quais
o agente de linguagem do \cref{cap:agente} calcula uma \emph{razão de oscilação}
explícita, e um dos motivos pelos quais o desempenho da classe 3 melhora
substancialmente quando se usa segmentação temporal (\cref{cap:segmentacao}) em vez de
classificação de janela agregada.

\subsection{Classe 4 --- Instabilidade de fluxo}

\textbf{Mecanismo.} Regime de escoamento não estacionário sem a periodicidade regular
da golfada severa. Pode originar-se de instabilidade de gás de elevação, de
acoplamento entre reservatório e coluna, ou de operação próxima ao limite de
estabilidade.

\textbf{Assinatura.} Variabilidade aumentada em todos os sinais, sem padrão dominante.

\textbf{Advertência.} A classe 4 é \emph{sintomática}, não causal. Instabilidade de
fluxo é aquilo que se observa; o que a causou pode ser qualquer coisa. Essa
característica explica por que ela apresenta o pior desempenho médio do livro --- 69\,\%
no ensemble binário (\cref{cap:closedworld}), 18,5\,\% na primeira escolha do agente
(\cref{cap:resultados-agente}) --- e por que foi deliberadamente excluída de dois dos
três estudos com modelos de linguagem.

\subsection{Classe 5 --- Perda rápida de produtividade}

\textbf{Mecanismo.} Queda súbita do índice de produtividade do poço, por dano à
formação, colapso de canhoneado, produção de finos ou bloqueio na região próxima ao
poço.

\textbf{Assinatura.} Queda de vazão com aumento da pressão de fundo (o reservatório
deixa de conseguir entregar fluido à mesma taxa).

\textbf{Advertência.} Como a classe 4, é sintomática. O \cref{cap:mahalanobis} reporta
para ela o pior resultado de detecção de novidade de toda a Tabela correspondente:
67,5\,\% no espaço latente, contra mais de 96\,\% para as demais classes.

\subsection{Classe 6 --- Restrição rápida no PCK}

\textbf{Mecanismo.} Obstrução súbita da válvula reguladora de produção, por detrito,
falha mecânica ou deposição abrupta.

\textbf{Assinatura.} Aumento da pressão a montante da \ing{choke} (P-TPT) com queda da
pressão a jusante. Contraste com a classe 2: ali a divergência ocorre \emph{através da
DHSV}, no fundo; aqui, \emph{através da choke}, na superfície.

\subsection{Classe 7 --- Incrustação no PCK}

\textbf{Mecanismo.} Deposição gradual de sais inorgânicos --- carbonato de cálcio,
sulfato de bário --- na passagem da válvula reguladora, reduzindo progressivamente a
área efetiva.

\textbf{Assinatura.} Mesma direção da classe 6, mas em escala de tempo de semanas a
meses. A distinção entre 6 e 7 é, essencialmente, \emph{uma distinção de derivada}.

\subsection{Classe 8 --- Hidrato na linha de produção}

\textbf{Mecanismo.} Hidratos de gás natural são estruturas cristalinas em que moléculas
de água formam gaiolas que aprisionam moléculas de metano, etano e propano. Formam-se
em condições de alta pressão e baixa temperatura --- exatamente as condições de uma
linha de produção submarina em águas profundas. O tampão resultante obstrui o
escoamento.

\textbf{Assinatura.} Aumento da pressão a montante do tampão, queda a jusante, e queda
de temperatura. A dinâmica é de horas.

\textbf{Destaque.} A classe 8 é o objeto do experimento central do
\cref{cap:closedworld}, em que é comparada com um modelo termodinâmico analítico, e é
a classe escondida no experimento de conjunto aberto. No \cref{cap:mahalanobis}, ela
apresenta a razão de separação mais extrema já observada nos experimentos deste livro
--- 186,5 vezes no espaço latente.

\subsection{Classe 9 --- Hidrato na linha de serviço}

\textbf{Mecanismo.} O mesmo da classe 8, mas na linha de serviço (anular ou linha de
injeção de gás), que opera em condições distintas e é menos instrumentada.

\textbf{Assinatura.} Semelhante à classe 8, porém observada indiretamente, por reflexo
sobre a linha de produção.

\textbf{Advertência.} A classe 9 é a mais escassa em número de instâncias e a de pior
desempenho consistente. No \cref{cap:resultados-agente}, o agente acerta 7,1\,\% na
primeira tentativa e não converge para um nome consensual quando solicitado a batizar
o evento como novidade.

\section{Além do 3W: eventos de completação inteligente}
\label{sec:dominio-icv}

Dois eventos discutidos neste livro não pertencem ao conjunto 3W. Vêm de dados
proprietários e cumprem um papel metodológico específico: \emph{servem de novidade
genuína}, isto é, de classe que o sistema comprovadamente nunca viu.

\begin{description}[leftmargin=0pt,style=nextline,itemsep=0.5em]

\item[Falha de ICV.] A válvula de controle de intervalo não responde ao comando, ou
responde parcialmente, ou apresenta vazamento entre zonas. A assinatura é um degrau de
pressão e vazão que não corresponde a nenhum comando registrado. No
\cref{cap:mundoaberto}, instâncias de falha de ICV são apresentadas a um sistema
treinado apenas nas nove classes do 3W: 48\,\% são classificadas como pertencentes a
alguma classe conhecida e 52\,\% são corretamente marcadas como desconhecidas e
encaminhadas ao estágio de agrupamento.

\item[Operação de válvula.] Manobras programadas de abertura e fechamento. São
\emph{eventos normais} cuja assinatura se assemelha à de anomalias --- o caso clássico
do falso positivo legítimo. Um sistema que não distinga uma operação programada de um
fechamento espúrio é inútil em campo.

\end{description}

\begin{destaque}[Por que isso importa metodologicamente]
Avaliar detecção de novidade escondendo uma classe do próprio conjunto de treinamento é
válido, mas otimista: a classe escondida foi gerada pelo mesmo processo de coleta,
rotulada pela mesma equipe, com a mesma instrumentação. Uma classe genuinamente externa
--- vinda de outro conjunto de poços, com outra completação --- é um teste mais duro.
Os dois eventos acima cumprem esse papel.
\end{destaque}

\section{Uma taxonomia por assinatura}

Encerramos com uma reorganização das nove classes que será útil em todos os capítulos
seguintes. Em vez de agrupá-las por mecanismo, agrupamo-las por \emph{aquilo que o
algoritmo vê}.

\begin{center}
\small
\begin{tabular}{p{3.4cm} p{2.0cm} p{7.2cm}}
\toprule
\textbf{Assinatura} & \textbf{Classes} & \textbf{Implicação para o método} \\
\midrule
Divergência entre dois sinais de pressão & 2, 6 & Facilmente capturada por
descritores de relação entre sensores; alto desempenho universal. \\[0.3em]
Oscilação periódica & 3 & Exige preservação da estrutura temporal; penaliza
agregação por média. \\[0.3em]
Rampa lenta & 1, 7 & Exige janela longa; confunde-se com deriva de sensor e com
declínio natural. \\[0.3em]
Degrau com queda de temperatura & 8, 9 & Bem separável quando há dados de
temperatura confiáveis; a classe 9 sofre por escassez. \\[0.3em]
Aumento de variabilidade sem padrão & 4, 5 & Sintomática; baixa separabilidade
intrínseca. Candidata a tratamento diferenciado. \\
\bottomrule
\end{tabular}
\end{center}

Essa tabela antecipa a proposta de \emph{divisão de trabalho por classe} do
\cref{cap:sintese}: em vez de exigir que um único modelo trate igualmente as cinco
famílias de assinatura, atribui-se a cada família o componente que a trata melhor.

\resumocapitulo{%
Um poço submarino é observado por um punhado de sensores --- essencialmente pressão e
temperatura no fundo (PDG) e na cabeça (TPT) --- cujas leituras estão sujeitas a
congelamento, perda de telemetria e deriva. As nove classes de anomalia do conjunto 3W
diferem radicalmente em mecanismo físico, escala temporal e assinatura: há degraus
nítidos (fechamento espúrio de DHSV), oscilações periódicas (golfadas), rampas lentas
(incrustação) e classes meramente sintomáticas (instabilidade de fluxo, perda de
produtividade). Essa heterogeneidade é a razão estrutural pela qual nenhum modelo único
apresenta bom desempenho em todas as classes, e é o ponto de partida da arquitetura
modular defendida neste livro. Eventos de completação inteligente --- falha de ICV e
operação de válvula --- fornecem, por serem externos ao 3W, testes rigorosos de
detecção de novidade.}

\begin{exercicios}
\item Desenhe, à mão, as curvas esperadas de P-PDG e P-TPT ao longo de trinta minutos
      para: (a) fechamento espúrio da DHSV; (b) restrição rápida no PCK. Indique
      claramente em que ponto as duas assinaturas se distinguem e qual descritor
      numérico capturaria essa distinção.

\item A diferença $P_{\text{PDG}} - P_{\text{TPT}}$ combina componente hidrostática e
      componente friccional. Para cada uma das classes 1, 2, 5 e 8, diga se você espera
      que essa diferença aumente, diminua ou permaneça estável, e justifique fisicamente.

\item As classes 6 e 7 compartilham a mesma assinatura qualitativa e diferem apenas na
      escala de tempo. Proponha dois descritores numéricos que as separariam de forma
      robusta e discuta como o comprimento da janela de observação afeta cada um deles.

\item Argumente a favor e contra a seguinte proposta: ``as classes 4 e 5 deveriam ser
      removidas do catálogo de classificação e tratadas apenas como sinalizadores de
      atenção''. Considere as consequências para a métrica global, para o operador e
      para o registro histórico de eventos.

\item Um sensor de fundo apresenta deriva de $+2$\,bar por mês. Descreva o efeito
      dessa deriva sobre: (a) um detector de anomalia baseado em limiar absoluto;
      (b) um detector baseado em erro de reconstrução de autoencoder treinado há seis
      meses; (c) um descritor de variação relativa à linha de base imediatamente
      anterior ao evento. Ordene as três abordagens por robustez.

\item Consulte a documentação do conjunto 3W \citep{VazVargas2026_3w2} e liste todas as
      variáveis disponíveis. Para cada variável não utilizada neste livro, argumente se
      ela poderia ou não contribuir para separar as classes 4 e 5.
\end{exercicios}

\begin{leituras}
\item \citet{VARGAS2019} descrevem a origem operacional de cada classe do 3W. É a fonte
      primária para o conteúdo da \cref{sec:dominio-taxonomia}.
\item \citet{VazVargas2026_3w2} apresentam a versão 2.0 do conjunto, com classes
      adicionais e revisão da rotulação.
\item \citet{Safamirzaei2015} tratam da termodinâmica de formação de hidratos, base do
      modelo analítico usado no \cref{cap:closedworld}.
\item \citet{aranha2023system} e \citet{Aranha2024} descrevem sistemas de monitoramento
      em operação, com ênfase nos requisitos práticos que raramente aparecem em
      artigos metodológicos.
\end{leituras}
